\begin{answer}

    \textbf{Number of Iterations} $\rightarrow$ The semi-supervised setting converged usually around in 1/3rd of the time ($\sim$35 Iterations), of the unsupervised setting ($\sim$123 Iterations).


    \textbf{Stability-wise} $\rightarrow$ Unsupervised setting is relatively unstable around the border of clusters, where certain points are sometimes part of different clusters. Looking at the loss, all three different trials have more or less the same loss but do have some small changes in the ballpark of $10^{-2}$. The semi-supervised setting seems to be even more stable though, where two of the trials even have the same log-likelihood, so they probably converged to the same minimum/maximum. In the plot-trials there is not any variation visible with the eye.

    
    \textbf{Quality} $\rightarrow$ Assuming we take the log-likelihood as the quality of the model, the semi-supervised models do perform better. And taking into account the information of how the datapoints got sampled, the resulting clusters of the semi-supervised setting do reflect 3 low-variance clusters and 1 high-variance cluster, whereas the unsupervised setting has two high-variance clusters (smaller variance than the single high-variance cluster from semi-supervised clustering).
\end{answer}
